\section{Web3 RegTech Ecosystem Analysis}
\label{sec:ecosystem}

This section addresses \textbf{RQ1} and \textbf{RQ2} through comprehensive analysis of the Web3 RegTech ecosystem. We examine how blockchain-native architectures enable novel compliance capabilities, analyze the current state of commercial platforms and academic research, and identify critical gaps limiting compliance effectiveness. Our analysis of 41 operational platforms and 28 research prototypes reveals both substantial progress and persistent challenges. Complete platform and prototype catalogs with detailed categorization and verification sources are provided in \autoref{tab:platform_catalog}.


\begin{table*}[htbp]
\centering
\caption{Web3 RegTech Platform Catalog. \textbf{Chain coverage levels}: L0 = chain-specific or limited single-chain coverage; L1 = limited multi-chain coverage; L2 = multi-chain coverage; L3 = broad multi-chain coverage; L4 = comprehensive multi-chain coverage; P = protocol-agnostic (travel-rule networks where chain coverage is not directly comparable).}
\label{tab:platform_catalog}
\tiny
\begin{tabularx}{\textwidth}{p{1.5cm}p{0.9cm}X p{1.1cm}p{1.2cm}p{1.1cm}X p{0.9cm}}
\hline
\textbf{Platform} & \textbf{Type} & \textbf{Primary Function} & \textbf{Architecture} & \textbf{Chain Coverage} & \textbf{Compliance Protocols} & \textbf{Technical Approach} & \textbf{Maturity} \\
\hline
\multicolumn{8}{l}{\textit{\textbf{Comprehensive General-Purpose Platforms (11)}}} \\
\hline
Chainalysis \footnote{\url{chainalysis.com}} & Web3 & Tx Monitoring, Forensics & SaaS & L4 & Address/Wallet Screening, KYT & Graph, Attribution & Enterprise \\
Elliptic \footnote{\url{elliptic.co}} & Web3 & Risk Scoring, Forensics & SaaS & L4 & Address/Wallet Screening, KYT & ML/DL, Graph & Enterprise \\
CipherTrace\footnote{Acquired by Mastercard 2021, continues operations} \footnote{\url{ciphertrace.com}} & Web3 & Tx Monitoring & SaaS & L4 & Address/Wallet Screening, KYT & Heuristic, ML & Enterprise \\
TRM Labs \footnote{\url{trmlabs.com}} & Web3 & Risk Scoring, Compliance & SaaS & L3 & Address/Wallet Screening, KYT & Graph, ML & Production \\
Merkle Science \footnote{\url{merklescience.com}} & Web3 & Tx Monitoring & SaaS & L3 & Address/Wallet Screening, KYT & Graph Analysis & Production \\
Scorechain \footnote{\url{scorechain.com}} & Web3 & Risk Scoring & SaaS & L2 & Address/Wallet Screening, KYT & Heuristic, ML & Production \\
Coinfirm\footnote{Acquired by Lukka May 2024, platform integrated into Lukka ecosystem} \footnote{\url{coinfirm.com}} & Web3 & Risk Scoring, AML & SaaS & L2 & Address/Wallet Screening, KYT & Rule-Based, Heuristic & Production \\
Crystal Intelligence \footnote{\url{crystalintelligence.com}} & Web3 & Forensics, Investigation & Hybrid & L2 & Address/Wallet Screening, KYT & Graph Analysis & Production \\
Solidus Labs \footnote{\url{soliduslabs.com}} & Web3 & Market Surveillance & SaaS & L3 & KYT & ML, Anomaly & Production \\
BIG\footnote{MRR +36.5\% in 2024; 2,500+ trained professionals. Source: BIGG Digital Assets 2024 Financial Results (Feb 2025). \url{https://biggdigitalassets.com}} \footnote{\url{blockchaingroup.io}} & Web3 & Forensics, Investigation & Hybrid & L3 & Address/Wallet Screening, KYT & Graph, Manual & Production \\
Chainalysis KYT \footnote{\url{chainalysis.com/solutions/kyt}} & Web3 & Enterprise Tx Monitoring & On-Premise & L4 & KYT & Graph, Real-Time & Enterprise \\
\hline
\multicolumn{8}{l}{\textit{\textbf{DeFi-Specialized Platforms (8)}}} \\
\hline
Forta Network \footnote{\url{forta.org}} & Web3 & Real-Time Monitoring & Decentralized & L4 & KYT & Rule-Based, On-Chain & Production \\
Lossless Protocol \footnote{\url{lossless.io}} & Web3 & Fraud Prevention & Decentralized & L1 & KYT & Heuristic, Governance & Production \\
Hypernative\footnote{Company claims detecting \$2.2B in losses. Source: \url{https://www.hypernative.io}} \footnote{\url{hypernative.io}} & Web3 & DeFi Monitoring, MEV & SaaS & L3 & KYT, Contract & ML, Real-Time & Production \\
Dedaub\footnote{Sponsored DeFi Security Summit 2024, Bangkok Nov 7-9} \footnote{\url{dedaub.com}} & Web3 & Smart Contract Security & SaaS & L0 & Contract Analysis & Semantic, Formal & Production \\
Chaos Labs \footnote{\url{chaoslabs.xyz}} & Web3 & DeFi Risk Management & SaaS & L2 & Protocol Risk, KYT & Simulation, Economic & Production \\
Gauntlet\footnote{Company claims monitoring \$42B+ in assets; Forbes Fintech 50 (2024). Source: \url{https://gauntlet.network}} \footnote{\url{gauntlet.network}} & Web3 & DeFi Risk Optimization & SaaS & L2 & Protocol Risk & Simulation, Agent-Based & Production \\
Nansen \footnote{\url{nansen.ai}} & Web3 & On-Chain Analytics & SaaS & L3 & KYA & Entity Labeling, Behavioral & Production \\
Arkham \footnote{\url{arkhamintelligence.com}} & Web3 & Entity Attribution & SaaS & L2 & KYA & Attribution, Community & Production \\
\hline
\multicolumn{8}{l}{\textit{\textbf{Cross-Chain Analytics Platforms (12)}}} \\
\hline
AnChain.AI \footnote{\url{anchain.ai}} & Web3 & Cross-Chain Investigation & SaaS & L3 & Address/Wallet Screening, KYT & AI/ML, Graph & Production \\
Bitquery \footnote{\url{bitquery.io}} & Web3 & Blockchain Data Infra & SaaS & L4 & KYT & Graph DB, Query & Production \\
Covalent \footnote{\url{covalenthq.com}} & Web3 & Unified Blockchain API & SaaS & L4 & KYT & Data Indexing & Production \\
Elementus\footnote{Company claims 650M+ addresses attributed; raised \$27M+. Source: \url{https://elementus.io}} \footnote{\url{elementus.io}} & Web3 & Cross-Chain Forensics & SaaS & L2 & Address/Wallet Screening, KYT & Graph, Attribution & Production \\
Bison Trails\footnote{Acquired by Coinbase, now operates as Coinbase Cloud} \footnote{\url{coinbase.com/cloud}} & Web3 & Multi-Chain Infrastructure & SaaS & L4 & KYT & Data Indexing & Production \\
Notabene \footnote{\url{notabene.id}} & Web3 & Travel Rule Platform & SaaS & P & KYC, KYB & Encrypted Messaging & Production \\
Sygna Bridge \footnote{\url{sygna.io}} & Web3 & Travel Rule Network & Hybrid & P & KYC, KYB & Registry, Crypto & Production \\
TRISA \footnote{\url{trisa.io}} & Web3 & Travel Rule Network & Decentralized & P & KYC, KYB & Federated, PKI & Production \\
21 Analytics \footnote{\url{https://www.21analytics.ch}} & Web3 & VASP Compliance & SaaS & L2 & KYC, Address/Wallet Screening, KYT & Travel Rule, Monitoring & Production \\
Shyft Network \footnote{\url{shyft.network}} & Web3 & Compliance Infrastructure & Hybrid & L1 & KYC, KYB, Travel & Attestation, Identity & Production \\
TransactID \footnote{\url{transactid.com}} & Web3 & Travel Rule Messaging & SaaS & P & KYC, KYB & Secure Messaging & Production \\
CipherBlade \footnote{\url{cipherblade.com}} & Web3 & Investigation, Recovery & Hybrid & L2 & KYA & Manual Investigation & Production \\
\hline
\multicolumn{8}{l}{\textit{\textbf{Privacy-Preserving Platforms (3)}}} \\
\hline
Privado ID \footnote{\url{privado.id}} & Web3 & Privacy-Preserving ID & Decentralized & L0 & KYC, KYA, Credentials & ZKP, SSI & Beta \\
Polygon ID \footnote{\url{polygon.technology/polygon-id}} & Web3 & Decentralized Identity & Decentralized & L0 & KYC, KYA, Credentials & ZKP, DID & Production \\
KILT Protocol \footnote{\url{kilt.io}} & Web3 & Credential Platform & Decentralized & L0 & KYC, KYB, Attestation & Credentials, Anchoring & Production \\
\hline
\multicolumn{8}{l}{\textit{\textbf{Traditional Platforms with Web3 Capabilities (7)}}} \\
\hline
ComplyAdvantage \footnote{\url{complyadvantage.com}} & Traditional & AML, Sanctions & SaaS & L1 & KYC, KYB, KYT & ML, Risk Scoring & Enterprise \\
Refinitiv World-Check \footnote{\url{refinitiv.com}} & Traditional & Sanctions, PEP & SaaS & L0 & KYC, KYB & Database, Screening & Enterprise \\
LexisNexis Bridger XG \footnote{\url{lexisnexis.com}} & Traditional & AML, Compliance & SaaS & L0 & KYC, KYB & Database, Network & Enterprise \\
NICE Actimize \footnote{\url{niceactimize.com}} & Traditional & Tx Monitoring & On-Premise & L0 & KYT & Rule-Based, ML & Enterprise \\
Fenergo \footnote{\url{fenergo.com}} & Traditional & Client Lifecycle Mgmt & SaaS & L0 & KYC, KYB & Workflow Automation & Enterprise \\
AMLBot \footnote{\url{amlbot.com}} & Traditional & Telegram Compliance & SaaS & L2 & Address/Wallet Screening, KYT & Address Screening & Production \\
ComplyAdvantage Crypto \footnote{\url{complyadvantage.com/crypto}} & Traditional & Crypto AML Extension & SaaS & L2 & KYT & Risk Database & Production \\
\hline
\end{tabularx}
\end{table*}



\subsection{Platform Architecture and Novel Capabilities}

Traditional RegTech and Web3 RegTech rest on fundamentally incompatible data architectures. Traditional platforms operate within centralized financial infrastructure where institutions maintain exclusive control over customer data and transaction records in proprietary databases accessible only through bilateral agreements. Customer identity verification relies on document-based processes where individuals submit government-issued identification verified through institutional processes. Transaction monitoring employs rule-based surveillance within single-institution data silos, fundamentally limiting analytical scope. Trust assumptions center on institutional accountability verified through periodic regulatory examinations.

Web3 RegTech exploits blockchain's transparency and immutability to enable compliance architectures impossible in centralized systems \cite{nakamoto2008bitcoinpeertopeerelectronic,narayanan2016bitcoincryptocurrencytechnologies}. Public blockchains provide universal read access to complete transaction histories, eliminating data sharing agreements as prerequisites for comprehensive monitoring. Any party can independently construct transaction graphs spanning entire blockchain networks, analyze fund flows through arbitrary transaction chains, identify behavioral patterns across pseudonymous addresses, and verify historical activities without requiring institutional cooperation \cite{bonneau2015sokresearchperspectives}. The challenge shifts from data acquisition to interpretation, transforming pseudonymous address interactions into meaningful compliance intelligence through attribution, clustering, and behavioral analysis.

\noindent\textbf{Comprehensive Transaction Graph Analysis.}
Unlike traditional systems where transaction visibility is limited to single-institution relationships, Web3 RegTech constructs transaction graphs spanning entire blockchain networks. Multi-hop fund tracing follows flows through arbitrary chain lengths, revealing ultimate sources and destinations despite intermediary addresses \cite{moser2013inquirymoneylaundering,ron2013quantitativeanalysisfull}. Entity clustering algorithms aggregate addresses into clusters representing common controlling entities through heuristic analysis \cite{meiklejohn2013fistfulbitcoinscharacterizing,reid2013analysisanonymitybitcoin,androulaki2013evaluatinguserprivacy}. The common input ownership heuristic exploits the observation that addresses used as inputs in the same Bitcoin transaction typically belong to the same wallet. Change address detection identifies outputs returning funds to senders based on behavioral patterns.

Network community detection identifies clusters of frequently interacting addresses potentially revealing organizational structures or criminal networks \cite{wu2021detectingmixingservices,drezewski2015applicationsocialnetwork}. Behavioral pattern recognition distinguishes characteristic activities: exchanges exhibit high transaction volumes with diverse counterparties, mixing services show rapid multi-hop forwarding with amount randomization, darknet markets demonstrate regular small transactions to common deposit addresses \cite{moser2018empiricalanalysistraceability}. Contemporary research advances graph analysis through deep learning. Graph neural networks learn representations from transaction graph topology enabling classification and anomaly detection at scale \cite{weber2019antimoneylaunderingbitcoin,azad2025machinelearningblockchain,ancelotti2024reviewblockchainapplication}. Temporal graph learning models dynamic patterns capturing evolving behaviors \cite{arnold2024insightscaveatsmining,luo2025optimizingblockchainanalysis}.

Smart contract analytics complement address-level tracing by surfacing vulnerable contracts, abusive permission settings, and protocol misconfigurations that frequently precipitate laundering incidents. Surveys of Ethereum attack vectors and benchmarking studies on vulnerability scanners inform RegTech feature engineering, enabling automated assessment of re-entrancy, integer overflow, access control, and oracle-manipulation risks that propagate through DeFi dependency graphs \cite{atzei2017surveyattacksethereum,baets2024vulnerabilitydetectionsmart}. Integrating these insights into entity profiles allows compliance teams to correlate suspicious fund movements with exploited contracts and proactively geofence exposure to high-risk bytecode patterns.

\noindent\textbf{Real-Time Risk Assessment.}
Public blockchain transparency enables risk assessment before or immediately upon transaction confirmation. Transactions broadcast to blockchain networks become visible in mempools before block inclusion, enabling pre-confirmation screening. Taint analysis quantifies the proportion of transaction inputs traceable to known illicit sources, providing interpretable risk scores \cite{moser2013inquirymoneylaundering,moser2014riskscoringbitcoin}. Address/wallet screening (often implemented as address risk scoring) aggregates multiple risk dimensions into composite scores reflecting historical behaviors, entity attributions, and network characteristics \cite{weber2019antimoneylaunderingbitcoin,pocher2023detectinganomalouscryptocurrency}. Transaction-level risk scoring evaluates individual transactions based on counterparty risk, amount anomalies, protocol interactions, and contextual factors. Modern approaches incorporate machine learning models trained on labeled datasets of illicit transactions, learning subtle patterns that rule-based systems miss \cite{alarab2020comparativeanalysisusing,lorenz2020machinelearningmethods,alotibi2022moneylaunderingdetection}.

Market-integrity analytics extend KYT beyond simple transfers by interpreting decentralized exchange order flow, wash-trading schemes, and NFT-based obfuscation. Compliance teams monitor routing bots, spoofed liquidity, and creator royalty evasion to identify manipulative behavior and off-ramp risks unique to tokenized markets \cite{victor2021detectingquantifyingwash,wang2021nonfungibletokennft}.

\noindent\textbf{Cross-Chain Analytics.}
Despite growing capabilities, comprehensive cross-chain fund tracking remains challenging. Bridge transaction monitoring identifies when assets transition between blockchains via cross-chain bridges, maintaining attribution across these transitions \cite{yousaf2019tracingtransactionscryptocurrency,zamyatin2021sokcommunicationdistributed,zhang2024securitycrosschainbridges}. Wrapped asset tracking follows tokens representing native assets from other chains. Cross-chain address clustering identifies addresses on different blockchains controlled by common entities based on behavioral correlations, timing patterns, and explicit linkages \cite{lin2024crosschainabnormaltransaction}. Unified entity graphs integrate activities across multiple networks creating comprehensive profiles of cross-chain behaviors \cite{coquide2025alamedaresearchmultitoken}.

\noindent\textbf{Smart Contract Interaction Analysis.}
Understanding DeFi protocol interactions provides compliance insights unavailable in traditional finance. Semantic transaction interpretation decodes smart contract function calls to extract meaningful semantics \cite{chen2020understandingethereumgraph,victor2019measuringethereumbasederc20}. Rather than presenting raw contract calls, semantic analysis translates these into human-readable descriptions. Contemporary approaches leverage language models to learn semantic mappings from transaction execution traces \cite{sun2025ethereumfrauddetection}. DeFi strategy recognition identifies complex multi-step operations including flash loans, arbitrage sequences, and yield farming strategies \cite{qin2021attackingdefiecosystem,daian2020flashboys20}. Protocol exposure tracking monitors which DeFi protocols customers interact with, assessing associated risks based on security audits, historical incidents, and governance quality \cite{sun2024sokcomprehensiveanalysis,ling2025sokstablecoindesigns}.

\noindent\textbf{Privacy-Preserving Compliance Verification.}
Cryptographic techniques including zero-knowledge proofs, verifiable credentials, and secure multi-party computation offer pathways for balancing regulatory requirements with privacy preservation \cite{ben-sasson2017scalablezeroknowledge,reed2022decentralizedidentifiersdids,zhang2024blockchainzeroknowledge}. While similar techniques exist in traditional finance contexts (e.g., private set intersection for fraud detection, secure multi-party computation for cross-institution analytics, and confidential computing for sensitive data processing), blockchain's public transparency and programmability enable more native integration of privacy-preserving compliance. However, commercial adoption in both Web3 and traditional systems remains minimal due to computational overhead, implementation complexity, regulatory acceptance uncertainty, and lack of standardized approaches. Discussion of privacy-preserving compliance architectures appears in \autoref{sec:discussion} research directions.

\subsection{Commercial Platform Landscape}

Analysis of 41 RegTech platforms indicates recent market expansion, with many Web3-native solutions in our sample launching since 2023 based on public announcements.\footnote{\url{https://www.coindesk.com/business/} and \url{https://www.theblock.co/} track blockchain compliance technology funding and launches.} This expansion appears correlated with increasing regulatory clarity, growing VASP compliance obligations, maturation of blockchain analytics techniques, and venture investment following high-profile enforcement actions.\footnote{\url{https://www.justice.gov/opa/pr/binance-and-ceo-plead-guilty-federal-charges} - Binance \$4.3B settlement (Nov 2023).}

\noindent\textbf{Market Evolution}

Early entrants focused on comprehensive general-purpose platforms providing transaction monitoring, risk scoring, and forensics across major blockchains. Platforms including Chainalysis, Elliptic, and CipherTrace established dominant positions through early mover advantages and extensive entity attribution databases built over years.\footnote{CipherTrace acquired by Mastercard 2021.} Recent entrants increasingly specialize in particular capabilities: DeFi-specific monitoring, cross-chain analytics, Travel Rule implementation, privacy-preserving compliance, and API-first integration platforms. This specialization suggests market maturation with established players dominating core capabilities while innovative entrants address emerging needs.

Chain coverage in our sample reflects prioritization based on network size and regulatory significance. Bitcoin and Ethereum show the broadest coverage. Major alternative Layer-1 blockchains including Binance Smart Chain, Polygon, Avalanche, Solana, and Tron show substantial but incomplete coverage. Ethereum Layer-2 solutions and emerging networks show lower coverage, representing compliance gaps.

\noindent\textbf{Capability Distribution}

Transaction monitoring appears near-universally as a core compliance requirement. Risk scoring is widespread, though methodology sophistication varies substantially. Address/wallet screening appears in 13 of 41 platforms, concentrated in comprehensive and cross-chain categories. Address attribution databases are common, with coverage quality representing a primary differentiator.\footnote{Entity attribution databases map blockchain addresses to real-world entities through clustering and exchange analysis.} Travel Rule support appears in roughly 15\% of platforms (6 of 41), reflecting ongoing protocol fragmentation.\footnote{\url{https://trisa.io/} - Travel Rule Information Sharing Architecture standardization effort; networks/solutions/standards include TRUST, Sygna Bridge, Notabene, and IVMS101.} DeFi analytics represent growing but not yet universal functionality.\footnote{\url{https://defillama.com/} tracks DeFi total value locked (TVL) exceeding \$100B across major protocols.} Cross-chain tracking appears as an advanced feature in a minority of platforms.\footnote{Major bridge hacks: Ronin (~624M USD, Mar 2022), Poly Network (~610M USD, Aug 2021), Wormhole (~320M USD, Feb 2022). Sources: Ronin Blog, Poly Network Medium, Wormhole Blog.} Privacy-preserving features show minimal adoption due to technical complexity and regulatory uncertainty.

\subsection{Academic Research Landscape}

Academic research explores novel compliance approaches often years ahead of commercial deployment. Graph neural networks enable deep learning on transaction graph structures, capturing complex patterns traditional heuristics miss \cite{weber2019antimoneylaunderingbitcoin,azad2025machinelearningblockchain,ancelotti2024reviewblockchainapplication,damico2025blockchainnetworkanalysis}. However, evaluations reveal that tree-based baselines often match GNN performance on benchmark datasets \cite{alarab2020comparativeanalysisusing,pourhabibi2020frauddetectionsystematic}, temporal generalization degrades as adversaries adapt evasion tactics \cite{monamo2016unsupervisedlearningrobust}, and label quality depends on law enforcement data and heuristic rules with inherent biases \cite{weber2019antimoneylaunderingbitcoin}.

Temporal graph learning models dynamic behavioral patterns \cite{arnold2024insightscaveatsmining,luo2025optimizingblockchainanalysis}. Money laundering detection research explores subgraph contrastive learning \cite{ouyang2024bitcoinmoneylaundering}, global-local graph attention \cite{li2025globallocalgraphattention}, dense flow analysis \cite{lin2024denseflowspottingcryptocurrency}, and early detection with path tracing \cite{cheng2023evolvepathtracer,cheng2024earlydetectionmalicious}. Phishing detection applies data augmentation and hybrid models \cite{chen2020phishingscamdetection,chen2024ethereumphishingscam,liu2024detectionethereumphishing}.

Privacy-preserving analytics research significantly exceeds commercial deployment. Zero-knowledge proofs for identity attribute verification enable proving specific claims without revealing credentials \cite{yang2024researchidentitydata,zhang2024blockchainzeroknowledge,ben-sasson2017scalablezeroknowledge}. Secure multi-party computation protocols enable collaborative analysis without data sharing \cite{bansod2022challengesmakingblockchain}. DeFi-specific analysis develops semantic interpretation techniques \cite{chen2020understandingethereumgraph,sun2025ethereumfrauddetection}, rug pull detection \cite{sun2024sokcomprehensiveanalysis}, and stablecoin risk assessment \cite{ling2025sokstablecoindesigns,financialactiontaskforcefatf2020reportg20socalled}. The synthesis in \autoref{tab:prototype_synthesis} summarizes research by theme (goals, methods, data, and limitations) rather than enumerating individual papers.


\begin{table*}[htbp]
\centering
\caption{Academic Research Prototype Synthesis}
\label{tab:prototype_synthesis}
\scriptsize
\begin{tabularx}{\textwidth}{XXXXX}
\hline
\textbf{Theme} & \textbf{Primary Goal} & \textbf{Common Technical Approaches} & \textbf{Typical Data / Labels} & \textbf{Validation \& Key Limitations} \\
\hline
Graph analysis \& clustering \cite{meiklejohn2013fistfulbitcoinscharacterizing,reid2013analysisanonymitybitcoin,ron2013quantitativeanalysisfull,androulaki2013evaluatinguserprivacy,wu2023tracerscalablegraphbased,lin2024crosschainabnormaltransaction} & Entity attribution, flow analysis, deanonymization & Heuristic clustering, graph statistics, GNNs & Full-chain graphs; heuristic entity labels; cross-chain linkages & Case studies, cluster precision, scalability; label bias and cross-chain gaps\\
Money laundering detection \cite{moser2013inquirymoneylaundering,moser2014riskscoringbitcoin,weber2019antimoneylaunderingbitcoin,ouyang2024bitcoinmoneylaundering,li2025globallocalgraphattention,lin2024denseflowspottingcryptocurrency} & Detect illicit flows and high-risk entities & Taint analysis, risk scoring, GNNs, contrastive learning, attention models & Labeled datasets (e.g., Elliptic), known mixers, synthetic labels & AUC/F1, temporal robustness; label scarcity and adversarial drift \\
Phishing \& fraud detection \cite{chen2020phishingscamdetection,chen2024ethereumphishingscam,liu2024detectionethereumphishing,pocher2023detectinganomalouscryptocurrency} & Identify scams, phishing, anomalies & Data augmentation, graph-based features, hybrid DL, unsupervised ML & Phishing labels, Etherscan reports, multi-chain activity & Precision/recall, early detection; noisy labels and limited transferability \\
Privacy protocol analysis \cite{kumar2017traceabilityanalysismoneros,kappos2018empiricalanalysisanonymity,wu2021detectingmixingservices} & Measure anonymity and leak risks & Heuristic/statistical deanonymization, behavioral detection & Privacy-coin chains, mixer datasets, known service addresses & Linkage accuracy and deanonymization rates; weak ground truth and ethical constraints \\
DeFi \cite{chen2020understandingethereumgraph,sun2025ethereumfrauddetection,qin2021attackingdefiecosystem,victor2021detectingquantifyingwash}  \& smart contract analysis & Interpret protocol behavior, attacks, market manipulation & Semantic decoding, graph construction, empirical taxonomy & Contract traces, protocol events, incident corpora & Coverage and case validation; fast-evolving protocols and semantic ambiguity \\
Privacy-preserving compliance \cite{yang2024researchidentitydata,zhang2024blockchainzeroknowledge,wang2024achievingprivacypreservingoptimizer} & KYC/AML with privacy guarantees & ZK credential proofs, secure computation, intent privacy & Prototype implementations, simulations & Proof size/latency and compliance guarantees; overhead and regulatory uncertainty \\
\hline
\end{tabularx}
\end{table*}

\subsection{Gap Analysis}

Systematic analysis reveals eight categories of disconnects between academic innovation and industry deployment, alongside six persistent challenges limiting current capabilities.

\subsubsection{Academia-Industry Gaps}

\noindent\textbf{Technology Readiness.} 
Many academic prototypes demonstrate algorithmic advances without addressing engineering challenges determining production viability. Scalability limitations appear frequently, with techniques achieving accuracy on small datasets failing to scale to billions of transactions. Real-time requirements demand sub-second response times, yet many approaches require minutes or hours for analysis. Fault tolerance and operational reliability receive limited attention, though production systems require comprehensive error handling.

\noindent\textbf{Regulatory Alignment.} 
Research sometimes addresses hypothetical scenarios disconnected from actual regulatory requirements. Privacy-preserving techniques enabling anonymous compliance verification face regulatory resistance from authorities requiring traditional identity verification. Decentralized compliance protocols eliminate identifiable accountable parties, conflicting with regulatory frameworks assuming hierarchical responsibility. Explainability requirements favor simpler heuristic approaches over black-box deep learning models despite performance advantages.

\noindent\textbf{Data Availability.} 
Research frequently assumes data availability that production environments cannot guarantee. Ground truth labels for illicit activities remain scarce. Off-chain intelligence including exchange attribution, dark web monitoring, and law enforcement sharing provides critical context that researchers cannot access. Cross-chain data integration requires comprehensive infrastructure beyond typical research scope.

\noindent\textbf{Adversarial Robustness.} 
Illicit actors actively adapt behaviors to evade detection through structuring, address rotation, and cross-chain migration \cite{albrecht2019usecryptocurrenciesmoney,guidara2022cryptocurrencymoneylaundering,wu2021detectingmixingservices}. Research evaluation against static datasets fails to capture this dynamic adversarial environment. Defensive techniques including adversarial training and ensemble methods remain nascent in production deployment.

\noindent\textbf{Cost-Benefit Considerations.} 
Research optimizes primarily for accuracy, treating computational costs as secondary. Production deployments must justify costs relative to operational benefits. Privacy-preserving cryptographic protocols imposing significant overhead may prove economically unviable despite technical elegance.

\noindent\textbf{Standardization and Interoperability.} 
Research prototypes implement custom interfaces without consideration for standardization. Production systems require standardized APIs, data formats, and integration protocols. The absence of compliance data standards creates fragmentation.

\noindent\textbf{Operational Workflow Integration.} 
Research focuses on technical detection without addressing operational workflows for alert management, investigation, case disposition, and regulatory reporting. Compliance teams require complete workflow systems, not isolated technical components.

\noindent\textbf{Validation and Benchmarking.} 
Lack of standardized benchmarks hinders comparison. Different research efforts use incompatible datasets, evaluation metrics, and experimental setups preventing direct comparison. Establishing shared benchmarks would substantially advance research quality and industry procurement decisions.

\subsubsection{Persistent Capability Challenges}

\noindent\textbf{Cross-Chain Tracking.} 
Cross-chain fund tracking faces significant challenges as assets transition between heterogeneous blockchain architectures. Attribution techniques including behavioral correlation, bridge monitoring, and wrapped asset tracking show promise \cite{yousaf2019tracingtransactionscryptocurrency,zamyatin2021sokcommunicationdistributed,zhang2024securitycrosschainbridges,li2025blockchaincrosschainbridge,lin2024crosschainabnormaltransaction}, though only 31\% of analyzed platforms report comprehensive multi-chain capabilities. Bridge protocol proliferation and architectural diversity create integration complexity.

\noindent\textbf{DeFi Interaction Analysis.}
Semantic interpretation of complex DeFi strategies including multi-protocol compositions, flash loan attacks, and MEV extraction remains incomplete. Intent-centric DeFi architectures further complicate interpretation by obscuring direct user intentions. In intent-centric systems, users express high-level goals while specialized solvers determine optimal execution paths across multiple protocols and chains \cite{mao2025motivationexecutiongrounded}. On-chain transactions reflect solver routing decisions rather than direct user choices, creating attribution challenges when determining whether specific counterparty interactions represent user intent or solver optimization. Current regulatory frameworks provide limited guidance on whether compliance obligations attach to users expressing intents, solvers executing them, or both.

\noindent\textbf{Privacy Protocol Analysis.} 
Privacy-enhancing technologies including mixing services, privacy coins, and confidential transactions effectively obscure fund flows. While some protocols exhibit structural vulnerabilities, many provide strong anonymity resisting current forensic approaches.

\noindent\textbf{Real-Time Scalability.} 
High-throughput networks and comprehensive multi-chain monitoring create scalability challenges. Sophisticated graph algorithms, machine learning inference, and semantic interpretation impose computational costs potentially prohibitive for real-time screening.

\noindent\textbf{Attribution Maintenance.} 
Entity attribution databases require continuous manual curation. The labor-intensive nature creates competitive moats favoring established platforms while limiting new entrant competitiveness.

\noindent\textbf{False Positive Management.} 
Transaction monitoring systems face substantial false positive rates, with vendors commonly reporting that many flagged transactions prove legitimate upon investigation \cite{alotibi2022moneylaunderingdetection}. Multi-tiered screening, feedback loop integration, and contextual enrichment attempt to balance detection sensitivity against operational feasibility, though standardized performance measurement remains absent.
